\documentclass[11pt,a4paper]{article}

%================================================
% PACKAGES
%================================================
\usepackage[utf8]{inputenc}
\usepackage[T1]{fontenc}
\usepackage[french]{babel}
\usepackage{lmodern}
\usepackage{amsmath,amssymb,amsthm}
\usepackage{geometry}
\usepackage{fancyhdr}
\usepackage{tikz}
\usepackage[most]{tcolorbox}
\usepackage{enumitem}
\usepackage{array}
\usepackage{tabularx}
\usepackage{xcolor}
\usepackage{multicol}
\usepackage{hyperref}

\geometry{margin=2.1cm}
\setlength{\parindent}{0pt}
\setlength{\parskip}{0.45em}

%================================================
% COULEURS
%================================================
\definecolor{bleuprepa}{RGB}{20,55,110}
\definecolor{vertprepa}{RGB}{20,105,70}
\definecolor{rougeprepa}{RGB}{160,35,35}
\definecolor{orangeprepa}{RGB}{180,100,20}
\definecolor{grisclair}{RGB}{245,247,250}

%================================================
% EN-TÊTE
%================================================
\pagestyle{fancy}
\fancyhf{}
\lhead{\textsc{Première générale -- Spécialité mathématiques}}
\rhead{\textsc{TD -- Suites numériques}}
\cfoot{\thepage}

%================================================
% BOÎTES
%================================================
\tcbset{
  enhanced,
  breakable,
  sharp corners=downhill,
  boxrule=0.7pt,
  colback=white,
  fonttitle=\bfseries
}

\newtcolorbox{methodbox}[1]{
  colframe=bleuprepa,
  colback=blue!3!white,
  title={Méthode -- #1}
}

\newtcolorbox{rememberbox}[1]{
  colframe=vertprepa,
  colback=green!3!white,
  title={À retenir -- #1}
}

\newtcolorbox{warningbox}[1]{
  colframe=rougeprepa,
  colback=red!3!white,
  title={Erreur classique -- #1}
}

\newtcolorbox{exobox}[2]{
  colframe=bleuprepa,
  colback=grisclair,
  title={Exercice #1 \hfill Difficulté : #2}
}

\newtcolorbox{solbox}[1]{
  colframe=vertprepa,
  colback=green!2!white,
  title={Correction de l'exercice #1}
}

%================================================
% COMMANDES
%================================================
\newcommand{\N}{\mathbb{N}}
\newcommand{\R}{\mathbb{R}}
\newcommand{\etoile}{\(\star\)}
\newcommand{\deuxetoiles}{\(\star\star\)}
\newcommand{\troisetoiles}{\(\star\star\star\)}
\newcommand{\quatreetoiles}{\(\star\star\star\star\)}
\newcommand{\cinqetoiles}{\(\star\star\star\star\star\)}

%================================================
% TITRE
%================================================
\begin{document}

\begin{center}
\begin{tikzpicture}
\node[
  fill=bleuprepa,
  text=white,
  rounded corners=3mm,
  inner sep=7mm,
  text width=0.92\linewidth,
  align=center
]{
  {\Huge \textbf{TD complet -- Suites numériques}}\\[2mm]
  {\Large Première générale -- Spécialité mathématiques}\\[1mm]
  {\large Poly d'entraînement progressif, rigoureux et corrigé}\\[1mm]
  {\small Niveau exigeant : applications, méthodes, démonstrations, problèmes longs}
};
\end{tikzpicture}
\end{center}

\vspace{0.5cm}

\begin{rememberbox}{Objectif du TD}
Ce TD vise à maîtriser les suites numériques au programme de Première spécialité mathématiques :
\begin{itemize}
  \item définition d'une suite par formule explicite ou par récurrence ;
  \item calcul de termes ;
  \item suites arithmétiques et géométriques ;
  \item étude du sens de variation ;
  \item somme des termes d'une suite géométrique ;
  \item modélisation de phénomènes discrets : évolution de population, placement financier, phénomènes répétés.
\end{itemize}
\end{rememberbox}

\begin{warningbox}{Consigne de rigueur}
Une suite n'est pas une fonction continue. Son indice est un entier naturel. On raisonne donc sur les termes \(u_n\), \(u_{n+1}\), \(u_{n+1}-u_n\), ou éventuellement \(\frac{u_{n+1}}{u_n}\) lorsque les termes sont strictement positifs.
\end{warningbox}

\tableofcontents
\newpage

%================================================
\section{Rappels de méthodes essentiels}

\begin{methodbox}{Calculer les termes d'une suite}
Pour calculer les premiers termes d'une suite :
\begin{itemize}
  \item si la suite est définie explicitement par \(u_n=f(n)\), on remplace \(n\) par \(0,1,2,\ldots\) ;
  \item si la suite est définie par récurrence, on part du terme initial puis on applique la relation de passage.
\end{itemize}
\end{methodbox}

\begin{methodbox}{Reconnaître une suite arithmétique}
Une suite \((u_n)\) est arithmétique lorsqu'il existe un réel \(r\) tel que, pour tout \(n\),
\[
u_{n+1}=u_n+r.
\]
Le nombre \(r\) s'appelle la raison. On a alors :
\[
u_n=u_0+nr.
\]
\end{methodbox}

\begin{methodbox}{Reconnaître une suite géométrique}
Une suite \((u_n)\) est géométrique lorsqu'il existe un réel \(q\) tel que, pour tout \(n\),
\[
u_{n+1}=q u_n.
\]
Le nombre \(q\) s'appelle la raison. On a alors :
\[
u_n=u_0 q^n.
\]
\end{methodbox}

\begin{methodbox}{Étudier les variations}
Pour étudier les variations de \((u_n)\), on peut :
\begin{itemize}
  \item étudier le signe de \(u_{n+1}-u_n\) ;
  \item si tous les termes sont strictement positifs, comparer \(\frac{u_{n+1}}{u_n}\) à \(1\).
\end{itemize}
\end{methodbox}

\begin{rememberbox}{Somme géométrique}
Si \(q\neq 1\), alors :
\[
1+q+q^2+\cdots+q^n=\frac{1-q^{n+1}}{1-q}.
\]
Plus généralement,
\[
u_0+u_1+\cdots+u_n=u_0\frac{1-q^{n+1}}{1-q}
\]
pour une suite géométrique de premier terme \(u_0\) et de raison \(q\neq 1\).
\end{rememberbox}

\newpage

%================================================
\section{Exercices d'application directe (APP)}

\begin{exobox}{1 -- Suite explicite}{\etoile}
On considère la suite définie pour tout \(n\in\N\) par
\[
u_n=3n^2-2n+5.
\]
\begin{enumerate}[label=\alph*)]
  \item Calculer \(u_0,u_1,u_2,u_3\).
  \item Calculer \(u_{10}\).
  \item Déterminer \(u_{n+1}\) en fonction de \(n\).
  \item Calculer \(u_{n+1}-u_n\), puis étudier le sens de variation de \((u_n)\).
\end{enumerate}
\end{exobox}

\begin{exobox}{2 -- Suite récurrente simple}{\etoile}
On définit la suite \((u_n)\) par
\[
u_0=4,\qquad u_{n+1}=2u_n-1.
\]
\begin{enumerate}[label=\alph*)]
  \item Calculer \(u_1,u_2,u_3,u_4\).
  \item La suite semble-t-elle croissante ?
  \item Calculer \(u_{n+1}-u_n\).
  \item Que faudrait-il savoir sur \(u_n\) pour conclure rigoureusement ?
\end{enumerate}
\end{exobox}

\begin{exobox}{3 -- Suite arithmétique donnée par récurrence}{\etoile}
Soit \((u_n)\) définie par
\[
u_0=7,\qquad u_{n+1}=u_n-3.
\]
\begin{enumerate}[label=\alph*)]
  \item Montrer que \((u_n)\) est arithmétique.
  \item Donner sa raison.
  \item Donner une formule explicite de \(u_n\).
  \item Calculer \(u_{25}\).
\end{enumerate}
\end{exobox}

\begin{exobox}{4 -- Suite géométrique donnée par récurrence}{\etoile}
Soit \((v_n)\) définie par
\[
v_0=5,\qquad v_{n+1}=1{,}2v_n.
\]
\begin{enumerate}[label=\alph*)]
  \item Montrer que \((v_n)\) est géométrique.
  \item Donner sa raison.
  \item Exprimer \(v_n\) en fonction de \(n\).
  \item Calculer \(v_6\), puis interpréter le résultat comme une hausse répétée de \(20\%\).
\end{enumerate}
\end{exobox}

\begin{exobox}{5 -- Identifier le type d'une suite}{\deuxetoiles}
Pour chacune des suites suivantes, dire si elle est arithmétique, géométrique ou ni l'une ni l'autre.
\[
a_n=4+5n,\qquad b_n=3\cdot 2^n,\qquad c_n=n^2+1,\qquad d_0=2,\ d_{n+1}=d_n+7.
\]
Justifier chaque réponse.
\end{exobox}

\begin{exobox}{6 -- Rang d'un dépassement}{\deuxetoiles}
Une suite arithmétique est définie par
\[
u_0=12,\qquad r=5.
\]
\begin{enumerate}[label=\alph*)]
  \item Donner \(u_n\).
  \item Résoudre l'inéquation \(u_n\geq 100\).
  \item Déterminer le premier rang à partir duquel \(u_n\geq 100\).
\end{enumerate}
\end{exobox}

\begin{exobox}{7 -- Placement avec intérêts composés}{\deuxetoiles}
On place \(800\) euros sur un compte rémunéré à \(3\%\) par an. On note \(C_n\) le capital après \(n\) années.
\begin{enumerate}[label=\alph*)]
  \item Donner \(C_0\).
  \item Exprimer \(C_{n+1}\) en fonction de \(C_n\).
  \item Identifier la nature de la suite \((C_n)\).
  \item Exprimer \(C_n\) en fonction de \(n\).
  \item Calculer le capital après \(10\) ans, à l'euro près.
\end{enumerate}
\end{exobox}

\begin{exobox}{8 -- Somme arithmétique}{\deuxetoiles}
Une suite arithmétique \((u_n)\) vérifie \(u_0=6\) et \(r=4\).
\begin{enumerate}[label=\alph*)]
  \item Calculer \(u_{20}\).
  \item Calculer \(S= u_0+u_1+\cdots+u_{20}\).
  \item Expliquer pourquoi on peut utiliser la formule :
  \[
  S=\frac{21(u_0+u_{20})}{2}.
  \]
\end{enumerate}
\end{exobox}

\begin{exobox}{9 -- Somme géométrique}{\deuxetoiles}
On considère la suite géométrique \((u_n)\) définie par
\[
u_0=2,\qquad q=3.
\]
\begin{enumerate}[label=\alph*)]
  \item Donner \(u_n\).
  \item Calculer \(S_n=u_0+u_1+\cdots+u_n\).
  \item Calculer \(S_5\).
\end{enumerate}
\end{exobox}

\begin{exobox}{10 -- Variation par différence}{\deuxetoiles}
On considère la suite définie par
\[
u_n=n^2-6n+8.
\]
\begin{enumerate}[label=\alph*)]
  \item Calculer \(u_{n+1}-u_n\).
  \item Déterminer pour quels entiers \(n\) on a \(u_{n+1}-u_n\geq 0\).
  \item En déduire le sens de variation de la suite.
\end{enumerate}
\end{exobox}

%================================================
\section{Exercices de méthode (METH)}

\begin{exobox}{11 -- Passage explicite/récurrent}{\deuxetoiles}
Soit \((u_n)\) définie par
\[
u_n=5-2n.
\]
\begin{enumerate}[label=\alph*)]
  \item Calculer \(u_{n+1}\).
  \item Montrer que \(u_{n+1}=u_n-2\).
  \item Donner une définition par récurrence de la suite.
  \item Préciser son sens de variation.
\end{enumerate}
\end{exobox}

\begin{exobox}{12 -- Suite géométrique et seuil}{\deuxetoiles}
On considère
\[
u_0=150,\qquad u_{n+1}=0{,}8u_n.
\]
\begin{enumerate}[label=\alph*)]
  \item Donner \(u_n\).
  \item Interpréter la raison \(0{,}8\) comme une baisse en pourcentage.
  \item Déterminer le plus petit entier \(n\) tel que \(u_n<50\).
\end{enumerate}
\end{exobox}

\begin{exobox}{13 -- Population avec augmentation constante}{\deuxetoiles}
Une ville compte \(12\,000\) habitants. Chaque année, sa population augmente de \(350\) habitants. On note \(P_n\) la population après \(n\) années.
\begin{enumerate}[label=\alph*)]
  \item Donner \(P_0\).
  \item Donner une relation de récurrence.
  \item Identifier la nature de la suite.
  \item Exprimer \(P_n\) en fonction de \(n\).
  \item Déterminer au bout de combien d'années la population dépassera \(20\,000\) habitants.
\end{enumerate}
\end{exobox}

\begin{exobox}{14 -- Population avec croissance en pourcentage}{\deuxetoiles}
Une population de bactéries vaut initialement \(5000\). Elle augmente de \(12\%\) par heure. On note \(B_n\) le nombre de bactéries après \(n\) heures.
\begin{enumerate}[label=\alph*)]
  \item Donner \(B_0\).
  \item Donner \(B_{n+1}\) en fonction de \(B_n\).
  \item Exprimer \(B_n\) en fonction de \(n\).
  \item Calculer \(B_8\), arrondi à l'unité.
\end{enumerate}
\end{exobox}

\begin{exobox}{15 -- Comparaison de deux placements}{\troisetoiles}
On compare deux placements.
\[
A_n=1000+60n,\qquad B_n=1000\cdot 1{,}04^n.
\]
\begin{enumerate}[label=\alph*)]
  \item Interpréter chaque placement.
  \item Calculer \(A_5\) et \(B_5\).
  \item Déterminer à l'aide d'un tableau le premier rang \(n\) tel que \(B_n>A_n\).
  \item Expliquer pourquoi le placement géométrique finit par devenir plus avantageux.
\end{enumerate}
\end{exobox}

\begin{exobox}{16 -- Suite auxiliaire}{\troisetoiles}
On considère la suite
\[
u_0=2,\qquad u_{n+1}=3u_n+4.
\]
On définit
\[
v_n=u_n+2.
\]
\begin{enumerate}[label=\alph*)]
  \item Calculer \(v_{n+1}\) en fonction de \(v_n\).
  \item Montrer que \((v_n)\) est géométrique.
  \item Exprimer \(v_n\), puis \(u_n\), en fonction de \(n\).
  \item Calculer \(u_5\).
\end{enumerate}
\end{exobox}

\begin{exobox}{17 -- Étude d'une suite quadratique}{\troisetoiles}
Soit
\[
u_n=n^2-4n+7.
\]
\begin{enumerate}[label=\alph*)]
  \item Calculer \(u_0,u_1,u_2,u_3,u_4\).
  \item Calculer \(u_{n+1}-u_n\).
  \item Étudier le signe de cette différence.
  \item En déduire le tableau de variations de \((u_n)\).
\end{enumerate}
\end{exobox}

\begin{exobox}{18 -- Somme des puissances de deux}{\troisetoiles}
On pose
\[
S_n=1+2+2^2+\cdots+2^n.
\]
\begin{enumerate}[label=\alph*)]
  \item Calculer \(S_0,S_1,S_2,S_3\).
  \item Conjecturer une formule explicite pour \(S_n\).
  \item Démontrer cette formule en utilisant la formule de la somme géométrique.
  \item Résoudre \(S_n\geq 1000\).
\end{enumerate}
\end{exobox}

\begin{exobox}{19 -- Suite décroissante positive}{\troisetoiles}
Soit
\[
u_n=\frac{5}{n+1}.
\]
\begin{enumerate}[label=\alph*)]
  \item Calculer \(u_0,u_1,u_2\).
  \item Montrer que \(u_{n+1}<u_n\) pour tout \(n\in\N\).
  \item La suite est-elle arithmétique ? géométrique ?
  \item Déterminer le premier rang \(n\) tel que \(u_n\leq 0{,}1\).
\end{enumerate}
\end{exobox}

\begin{exobox}{20 -- Algorithmique et suites}{\troisetoiles}
On considère l'algorithme suivant :
\[
\begin{array}{l}
u\leftarrow 3\\
\text{Pour } k \text{ allant de } 1 \text{ à } 6\\
\quad u\leftarrow 2u+1\\
\text{Fin Pour}
\end{array}
\]
\begin{enumerate}[label=\alph*)]
  \item Écrire les six valeurs successives prises par \(u\).
  \item Identifier la suite calculée.
  \item Définir cette suite par récurrence.
  \item Montrer que la suite auxiliaire \(v_n=u_n+1\) est géométrique.
\end{enumerate}
\end{exobox}

%================================================
\section{Exercices de démonstration (DEM)}

\begin{exobox}{21 -- Démonstration de la formule arithmétique}{\troisetoiles}
Soit \((u_n)\) une suite arithmétique de premier terme \(u_0\) et de raison \(r\).
\begin{enumerate}[label=\alph*)]
  \item Écrire \(u_1,u_2,u_3\) en fonction de \(u_0\) et \(r\).
  \item Conjecturer la formule générale.
  \item Démontrer que pour tout \(n\in\N\),
  \[
  u_n=u_0+nr.
  \]
  \item Appliquer avec \(u_0=-2\) et \(r=5\).
\end{enumerate}
\end{exobox}

\begin{exobox}{22 -- Démonstration de la formule géométrique}{\troisetoiles}
Soit \((u_n)\) une suite géométrique de premier terme \(u_0\) et de raison \(q\).
\begin{enumerate}[label=\alph*)]
  \item Écrire \(u_1,u_2,u_3\).
  \item Conjecturer une formule explicite.
  \item Démontrer que
  \[
  u_n=u_0q^n.
  \]
  \item Appliquer avec \(u_0=4\) et \(q=-2\).
\end{enumerate}
\end{exobox}

\begin{exobox}{23 -- Somme géométrique}{\troisetoiles}
Soit \(q\neq 1\). On pose
\[
S_n=1+q+q^2+\cdots+q^n.
\]
\begin{enumerate}[label=\alph*)]
  \item Calculer \(qS_n\).
  \item Calculer \(S_n-qS_n\).
  \item En déduire la formule de \(S_n\).
  \item Appliquer avec \(q=\frac12\) et \(n=5\).
\end{enumerate}
\end{exobox}

\begin{exobox}{24 -- Variations d'une suite arithmétique}{\troisetoiles}
Soit \((u_n)\) arithmétique de raison \(r\).
\begin{enumerate}[label=\alph*)]
  \item Exprimer \(u_{n+1}-u_n\).
  \item Montrer que si \(r>0\), alors la suite est strictement croissante.
  \item Montrer que si \(r<0\), alors la suite est strictement décroissante.
  \item Que se passe-t-il si \(r=0\) ?
\end{enumerate}
\end{exobox}

\begin{exobox}{25 -- Variations d'une suite géométrique positive}{\troisetoiles}
Soit \((u_n)\) une suite géométrique de premier terme \(u_0>0\) et de raison \(q>0\).
\begin{enumerate}[label=\alph*)]
  \item Exprimer \(u_{n+1}-u_n\).
  \item Montrer que si \(q>1\), la suite est strictement croissante.
  \item Montrer que si \(0<q<1\), la suite est strictement décroissante.
  \item Illustrer par deux exemples numériques.
\end{enumerate}
\end{exobox}

\begin{exobox}{26 -- Suite affine transformée}{\quatreetoiles}
On considère
\[
u_0=1,\qquad u_{n+1}=2u_n+3.
\]
\begin{enumerate}[label=\alph*)]
  \item Chercher un réel \(\alpha\) tel que la suite \(v_n=u_n+\alpha\) soit géométrique.
  \item Déterminer \(\alpha\).
  \item Exprimer \(u_n\) en fonction de \(n\).
  \item Vérifier la formule pour \(n=0,1,2\).
\end{enumerate}
\end{exobox}

\begin{exobox}{27 -- Inégalité par récurrence intuitive}{\quatreetoiles}
On considère
\[
u_0=1,\qquad u_{n+1}=u_n+2n+1.
\]
\begin{enumerate}[label=\alph*)]
  \item Calculer \(u_0,u_1,u_2,u_3,u_4\).
  \item Conjecturer une formule pour \(u_n\).
  \item Vérifier que cette formule est compatible avec la relation de récurrence.
  \item En déduire une formule explicite pour \(u_n\).
\end{enumerate}
\end{exobox}

\begin{exobox}{28 -- Monotonie d'une suite rationnelle}{\quatreetoiles}
Soit
\[
u_n=\frac{2n+3}{n+2}.
\]
\begin{enumerate}[label=\alph*)]
  \item Calculer \(u_0,u_1,u_2\).
  \item Calculer \(u_{n+1}-u_n\).
  \item Étudier le signe de cette différence.
  \item En déduire le sens de variation de \((u_n)\).
\end{enumerate}
\end{exobox}

\begin{exobox}{29 -- Encadrement}{\quatreetoiles}
Soit
\[
u_n=3+\frac{2}{n+1}.
\]
\begin{enumerate}[label=\alph*)]
  \item Montrer que pour tout \(n\in\N\), \(u_n>3\).
  \item Montrer que \(u_n\leq 5\).
  \item Étudier le sens de variation.
  \item Interpréter graphiquement l'idée d'une suite qui se rapproche de \(3\) sans jamais l'atteindre.
\end{enumerate}
\end{exobox}

\begin{exobox}{30 -- Somme pondérée}{\quatreetoiles}
Soit \((u_n)\) géométrique de premier terme \(u_0=6\) et de raison \(q=\frac13\).
\begin{enumerate}[label=\alph*)]
  \item Exprimer \(u_n\).
  \item Calculer \(S_n=u_0+\cdots+u_n\).
  \item Montrer que \(S_n<9\) pour tout \(n\).
  \item Expliquer pourquoi les termes ajoutés deviennent de plus en plus petits.
\end{enumerate}
\end{exobox}

%================================================
\section{Exercices de réflexion et de préparation approfondie (PREP)}

\begin{exobox}{31 -- Comparaison entre arithmétique et géométrique}{\quatreetoiles}
On considère
\[
u_n=10+5n,\qquad v_n=10\cdot 1{,}2^n.
\]
\begin{enumerate}[label=\alph*)]
  \item Calculer les cinq premiers termes de chaque suite.
  \item Déterminer expérimentalement le premier rang pour lequel \(v_n>u_n\).
  \item Justifier que \((u_n)\) croît par ajout constant et que \((v_n)\) croît par multiplication constante.
  \item Expliquer pourquoi une croissance géométrique peut dépasser une croissance arithmétique.
\end{enumerate}
\end{exobox}

\begin{exobox}{32 -- Modèle de population avec pertes}{\quatreetoiles}
Une population initiale de \(10\,000\) individus diminue chaque année de \(4\%\), mais reçoit aussi \(200\) nouveaux individus par an. On note \(P_n\) la population après \(n\) années.
\begin{enumerate}[label=\alph*)]
  \item Écrire la relation de récurrence.
  \item Calculer \(P_1,P_2,P_3\).
  \item Cette suite est-elle arithmétique ? géométrique ?
  \item On pose \(Q_n=P_n-5000\). Montrer que \((Q_n)\) est géométrique.
  \item En déduire \(P_n\).
\end{enumerate}
\end{exobox}

\begin{exobox}{33 -- Remboursement d'une dette}{\quatreetoiles}
Une personne emprunte \(5000\) euros. Chaque mois, la dette augmente de \(1\%\), puis la personne rembourse \(150\) euros. On note \(D_n\) la dette après \(n\) mois.
\begin{enumerate}[label=\alph*)]
  \item Écrire \(D_{n+1}\) en fonction de \(D_n\).
  \item Calculer \(D_1,D_2,D_3\).
  \item On pose \(E_n=D_n-15000\). Montrer que \((E_n)\) est géométrique.
  \item Exprimer \(D_n\).
  \item Déterminer si la dette peut s'annuler dans ce modèle.
\end{enumerate}
\end{exobox}

\begin{exobox}{34 -- Suite définie par une moyenne}{\quatreetoiles}
Soit
\[
u_0=10,\qquad u_{n+1}=\frac{u_n+4}{2}.
\]
\begin{enumerate}[label=\alph*)]
  \item Calculer \(u_1,u_2,u_3\).
  \item Montrer que si \(u_n>4\), alors \(u_{n+1}>4\).
  \item Montrer que \(u_{n+1}-u_n=\frac{4-u_n}{2}\).
  \item En déduire le sens de variation.
  \item Poser \(v_n=u_n-4\) et déterminer une formule explicite.
\end{enumerate}
\end{exobox}

\begin{exobox}{35 -- Phénomène répété}{\quatreetoiles}
Un médicament est injecté dans le sang. Chaque heure, \(30\%\) de la quantité présente est éliminée. On injecte ensuite \(5\) mg supplémentaires. On note \(M_n\) la quantité de médicament après \(n\) heures, juste après injection.
\begin{enumerate}[label=\alph*)]
  \item Donner une relation de récurrence.
  \item Calculer \(M_1,M_2,M_3\) si \(M_0=5\).
  \item Chercher une valeur stable \(L\) vérifiant \(L=0{,}7L+5\).
  \item Poser \(V_n=M_n-L\) et montrer que \((V_n)\) est géométrique.
  \item Exprimer \(M_n\).
\end{enumerate}
\end{exobox}

\begin{exobox}{36 -- Suite et inéquation exponentielle discrète}{\quatreetoiles}
On considère
\[
u_n=2\cdot 1{,}15^n.
\]
\begin{enumerate}[label=\alph*)]
  \item Interpréter la suite comme une augmentation de \(15\%\).
  \item Déterminer par essais successifs le premier rang \(n\) tel que \(u_n\geq 10\).
  \item Écrire un algorithme permettant de trouver ce rang.
  \item Justifier que ce rang existe.
\end{enumerate}
\end{exobox}

\begin{exobox}{37 -- Recherche d'une raison}{\quatreetoiles}
Une suite géométrique vérifie
\[
u_2=18,\qquad u_5=486.
\]
\begin{enumerate}[label=\alph*)]
  \item Exprimer \(u_5\) en fonction de \(u_2\) et de la raison \(q\).
  \item Déterminer \(q^3\).
  \item Trouver les raisons possibles si l'on travaille dans \(\R\).
  \item Donner une expression de \(u_n\) dans chacun des cas.
\end{enumerate}
\end{exobox}

\begin{exobox}{38 -- Suite arithmétique à partir de deux termes}{\quatreetoiles}
Une suite arithmétique vérifie
\[
u_4=17,\qquad u_{12}=49.
\]
\begin{enumerate}[label=\alph*)]
  \item Exprimer \(u_{12}-u_4\) en fonction de la raison \(r\).
  \item Déterminer \(r\).
  \item Déterminer \(u_0\).
  \item Exprimer \(u_n\).
\end{enumerate}
\end{exobox}

\begin{exobox}{39 -- Somme et équation}{\quatreetoiles}
On considère la suite géométrique \(u_n=3\cdot 2^n\).
\begin{enumerate}[label=\alph*)]
  \item Calculer \(S_n=u_0+\cdots+u_n\).
  \item Résoudre \(S_n=381\).
  \item Résoudre \(S_n\geq 1000\).
\end{enumerate}
\end{exobox}

\begin{exobox}{40 -- Étude complète d'un modèle discret}{\quatreetoiles}
Une entreprise produit \(1000\) objets le premier mois. Chaque mois, sa production augmente de \(8\%\), mais elle perd ensuite \(30\) objets à cause de défauts de production. On note \(P_n\) la production nette du mois \(n\).
\begin{enumerate}[label=\alph*)]
  \item Écrire une relation de récurrence.
  \item Montrer que cette suite n'est ni arithmétique ni géométrique.
  \item Déterminer une constante \(L\) telle que \(P_{n+1}-L=1{,}08(P_n-L)\).
  \item En déduire une formule explicite de \(P_n\).
  \item Étudier le sens de variation.
\end{enumerate}
\end{exobox}

%================================================
\section{Exercices type oraux exigeants}

\begin{exobox}{41 -- Oral type CCP : suite affine}{\quatreetoiles}
On considère
\[
u_0=0,\qquad u_{n+1}=0{,}6u_n+8.
\]
\begin{enumerate}[label=\alph*)]
  \item Calculer les quatre premiers termes.
  \item Trouver la valeur \(L\) vérifiant \(L=0{,}6L+8\).
  \item Montrer que \(v_n=u_n-L\) est géométrique.
  \item Exprimer \(u_n\).
  \item Déterminer le plus petit \(n\) tel que \(u_n>19\).
\end{enumerate}
\end{exobox}

\begin{exobox}{42 -- Oral type CCP : somme et modélisation}{\quatreetoiles}
Une personne verse \(100\) euros au début de chaque mois sur un compte rémunéré à \(0{,}5\%\) par mois. On note \(C_n\) le capital après \(n\) mois, juste après le versement.
\begin{enumerate}[label=\alph*)]
  \item Écrire une relation de récurrence.
  \item Calculer \(C_0,C_1,C_2\).
  \item Exprimer \(C_n\) à l'aide d'une somme géométrique.
  \item Simplifier cette expression.
  \item Calculer le capital après \(24\) mois.
\end{enumerate}
\end{exobox}

\begin{exobox}{43 -- Oral type Mines : seuil double}{\cinqetoiles}
On considère deux suites :
\[
u_n=500+30n,\qquad v_n=200\cdot 1{,}08^n.
\]
\begin{enumerate}[label=\alph*)]
  \item Interpréter les deux modèles.
  \item Déterminer expérimentalement le premier rang \(n\) tel que \(v_n>u_n\).
  \item Montrer que \((v_n)\) est croissante.
  \item Expliquer pourquoi une comparaison exacte est délicate en Première.
  \item Proposer un algorithme de recherche du rang.
\end{enumerate}
\end{exobox}

\begin{exobox}{44 -- Oral type Centrale : suite avec paramètre}{\cinqetoiles}
Soit \(a\) un réel positif. On définit
\[
u_0=a,\qquad u_{n+1}=2u_n+1.
\]
\begin{enumerate}[label=\alph*)]
  \item Calculer \(u_1,u_2,u_3\) en fonction de \(a\).
  \item Poser \(v_n=u_n+1\). Montrer que \((v_n)\) est géométrique.
  \item Exprimer \(u_n\).
  \item Déterminer pour quelles valeurs de \(a\) on a \(u_5>100\).
\end{enumerate}
\end{exobox}

\begin{exobox}{45 -- Oral type Mines : suite et encadrement}{\cinqetoiles}
Soit
\[
u_0=12,\qquad u_{n+1}=\frac{u_n+6}{2}.
\]
\begin{enumerate}[label=\alph*)]
  \item Calculer \(u_1,u_2,u_3\).
  \item Montrer que \(u_n>6\) pour tout \(n\).
  \item Montrer que \((u_n)\) est décroissante.
  \item Déterminer une expression explicite de \(u_n\).
  \item Trouver le premier rang tel que \(u_n<6{,}1\).
\end{enumerate}
\end{exobox}

\begin{exobox}{46 -- Oral type Centrale : suite géométrique alternée}{\cinqetoiles}
On considère
\[
u_n=3(-2)^n.
\]
\begin{enumerate}[label=\alph*)]
  \item Calculer les six premiers termes.
  \item La suite est-elle croissante ? décroissante ?
  \item Montrer qu'elle est géométrique.
  \item Calculer \(S_n=u_0+\cdots+u_n\).
  \item Discuter le signe de \(u_n\) selon la parité de \(n\).
\end{enumerate}
\end{exobox}

\begin{exobox}{47 -- Oral type CCP : modèle d'abonnement}{\quatreetoiles}
Un service en ligne possède \(2000\) abonnés. Chaque mois, \(5\%\) des abonnés se désabonnent, puis \(180\) nouveaux abonnés arrivent.
\begin{enumerate}[label=\alph*)]
  \item Définir une suite modélisant la situation.
  \item Calculer les trois premiers termes.
  \item Déterminer une valeur stable \(L\).
  \item Obtenir une formule explicite.
  \item Interpréter le résultat à long terme.
\end{enumerate}
\end{exobox}

\begin{exobox}{48 -- Oral type Mines : somme d'épargne}{\cinqetoiles}
Une personne économise \(50\) euros la première semaine, puis augmente son épargne de \(5\) euros chaque semaine.
\begin{enumerate}[label=\alph*)]
  \item Modéliser l'épargne hebdomadaire par une suite.
  \item Calculer la somme économisée pendant \(20\) semaines.
  \item Déterminer le nombre minimal de semaines nécessaires pour économiser au moins \(2000\) euros.
  \item Proposer une méthode algorithmique.
\end{enumerate}
\end{exobox}

\begin{exobox}{49 -- Oral type Centrale : formule cachée}{\cinqetoiles}
On considère
\[
u_0=2,\qquad u_{n+1}=u_n+3\cdot 2^n.
\]
\begin{enumerate}[label=\alph*)]
  \item Calculer \(u_1,u_2,u_3,u_4\).
  \item Écrire \(u_n\) comme une somme.
  \item Utiliser la somme géométrique pour obtenir une formule explicite.
  \item Vérifier la formule pour \(n=0,1,2\).
\end{enumerate}
\end{exobox}

\begin{exobox}{50 -- Grand exercice final type concours adapté Première}{\cinqetoiles}
On considère une suite \((u_n)\) définie par
\[
u_0=100,\qquad u_{n+1}=0{,}9u_n+12.
\]
\begin{enumerate}[label=\alph*)]
  \item Interpréter cette suite dans un contexte de stock : perte de \(10\%\), puis ajout de \(12\) unités.
  \item Calculer \(u_1,u_2,u_3\).
  \item Trouver la valeur stable \(L\) vérifiant \(L=0{,}9L+12\).
  \item Poser \(v_n=u_n-L\) et montrer que \((v_n)\) est géométrique.
  \item Exprimer \(u_n\) en fonction de \(n\).
  \item Étudier le sens de variation de \((u_n)\).
  \item Déterminer le premier rang \(n\) tel que \(u_n>119\).
\end{enumerate}
\end{exobox}

\newpage

%================================================
\section{Corrigé complet}

\begin{solbox}{1}
\[
u_0=5,\quad u_1=3-2+5=6,\quad u_2=12-4+5=13,\quad u_3=27-6+5=26.
\]
\[
u_{10}=3\cdot 100-20+5=285.
\]
Puis
\[
u_{n+1}=3(n+1)^2-2(n+1)+5=3n^2+4n+6.
\]
Donc
\[
u_{n+1}-u_n=(3n^2+4n+6)-(3n^2-2n+5)=6n+1.
\]
Pour tout \(n\in\N\), \(6n+1>0\). La suite est donc strictement croissante.
\end{solbox}

\begin{solbox}{2}
\[
u_1=2\cdot4-1=7,\quad u_2=13,\quad u_3=25,\quad u_4=49.
\]
La suite semble croissante. On a :
\[
u_{n+1}-u_n=2u_n-1-u_n=u_n-1.
\]
Pour conclure que la suite est croissante, il faut savoir que \(u_n\geq 1\) pour tout \(n\). Ici, les premiers termes le suggèrent et la relation \(u_{n+1}=2u_n-1\) conserve la propriété \(u_n\geq1\).
\end{solbox}

\begin{solbox}{3}
La relation \(u_{n+1}=u_n-3\) est de la forme \(u_{n+1}=u_n+r\) avec \(r=-3\). La suite est arithmétique de raison \(-3\). Donc :
\[
u_n=u_0+nr=7-3n.
\]
Ainsi :
\[
u_{25}=7-75=-68.
\]
\end{solbox}

\begin{solbox}{4}
On a \(v_{n+1}=1{,}2v_n\). La suite est géométrique de raison \(q=1{,}2\) et de premier terme \(v_0=5\). Donc :
\[
v_n=5\cdot 1{,}2^n.
\]
Puis
\[
v_6=5\cdot1{,}2^6\approx14{,}93.
\]
Multiplier chaque année par \(1{,}2\) correspond à une hausse de \(20\%\).
\end{solbox}

\begin{solbox}{5}
\[
a_n=4+5n
\]
est arithmétique de raison \(5\).  
\[
b_n=3\cdot 2^n
\]
est géométrique de raison \(2\).  
\[
c_n=n^2+1
\]
n'est ni arithmétique ni géométrique : les différences ne sont pas constantes et les quotients non plus.  
Enfin, \(d_{n+1}=d_n+7\) définit une suite arithmétique de raison \(7\).
\end{solbox}

\begin{solbox}{6}
La suite est arithmétique :
\[
u_n=12+5n.
\]
On résout :
\[
12+5n\geq100
\]
soit
\[
5n\geq88,
\]
donc
\[
n\geq17{,}6.
\]
Comme \(n\) est entier, le premier rang est \(n=18\).
\end{solbox}

\begin{solbox}{7}
On a \(C_0=800\). Une hausse de \(3\%\) correspond à une multiplication par \(1{,}03\), donc :
\[
C_{n+1}=1{,}03C_n.
\]
La suite est géométrique de raison \(1{,}03\). Donc :
\[
C_n=800\cdot1{,}03^n.
\]
Après \(10\) ans :
\[
C_{10}=800\cdot1{,}03^{10}\approx1075.
\]
Le capital est donc environ \(1075\) euros.
\end{solbox}

\begin{solbox}{8}
\[
u_{20}=6+20\cdot4=86.
\]
Il y a \(21\) termes de \(u_0\) à \(u_{20}\). Pour une suite arithmétique :
\[
S=\frac{\text{nombre de termes}\times(\text{premier}+\text{dernier})}{2}.
\]
Ainsi :
\[
S=\frac{21(6+86)}2=21\cdot46=966.
\]
\end{solbox}

\begin{solbox}{9}
\[
u_n=2\cdot3^n.
\]
Donc :
\[
S_n=2(1+3+\cdots+3^n)=2\frac{1-3^{n+1}}{1-3}=3^{n+1}-1.
\]
Ainsi :
\[
S_5=3^6-1=729-1=728.
\]
\end{solbox}

\begin{solbox}{10}
\[
u_{n+1}=(n+1)^2-6(n+1)+8=n^2-4n+3.
\]
Donc :
\[
u_{n+1}-u_n=(n^2-4n+3)-(n^2-6n+8)=2n-5.
\]
On a \(2n-5\geq0\) si \(n\geq3\). La suite décroît donc jusqu'au rang \(3\), puis croît à partir du rang \(3\).
\end{solbox}

\begin{solbox}{11}
\[
u_{n+1}=5-2(n+1)=5-2n-2=3-2n.
\]
Or
\[
u_n-2=(5-2n)-2=3-2n.
\]
Donc :
\[
u_{n+1}=u_n-2.
\]
Une définition par récurrence est :
\[
u_0=5,\qquad u_{n+1}=u_n-2.
\]
La suite est arithmétique de raison \(-2\), donc strictement décroissante.
\end{solbox}

\begin{solbox}{12}
\[
u_n=150\cdot0{,}8^n.
\]
Multiplier par \(0{,}8\), c'est conserver \(80\%\), donc subir une baisse de \(20\%\).  
On cherche :
\[
150\cdot0{,}8^n<50.
\]
Par essais :
\[
u_4=150\cdot0{,}8^4=61{,}44,\qquad u_5=49{,}152.
\]
Le premier rang est donc \(n=5\).
\end{solbox}

\begin{solbox}{13}
\[
P_0=12000,\qquad P_{n+1}=P_n+350.
\]
La suite est arithmétique de raison \(350\). Donc :
\[
P_n=12000+350n.
\]
On cherche :
\[
12000+350n>20000.
\]
Donc :
\[
350n>8000,
\]
et
\[
n>22{,}857.
\]
Le premier entier convenable est \(n=23\).
\end{solbox}

\begin{solbox}{14}
\[
B_0=5000,\qquad B_{n+1}=1{,}12B_n.
\]
La suite est géométrique de raison \(1{,}12\). Donc :
\[
B_n=5000\cdot1{,}12^n.
\]
Ainsi :
\[
B_8=5000\cdot1{,}12^8\approx12380.
\]
\end{solbox}

\begin{solbox}{15}
\(A_n\) représente un placement avec gain fixe de \(60\) euros par an.  
\(B_n\) représente un placement avec hausse annuelle de \(4\%\).  
\[
A_5=1000+300=1300,
\]
\[
B_5=1000\cdot1{,}04^5\approx1217.
\]
En testant :
\[
A_{15}=1900,\quad B_{15}\approx1801,
\]
\[
A_{20}=2200,\quad B_{20}\approx2191,
\]
\[
A_{21}=2260,\quad B_{21}\approx2279.
\]
Le premier rang est \(n=21\). Une croissance géométrique finit par dépasser une croissance arithmétique car elle applique un pourcentage à une base de plus en plus grande.
\end{solbox}

\begin{solbox}{16}
On a :
\[
v_{n+1}=u_{n+1}+2=3u_n+4+2=3u_n+6=3(u_n+2)=3v_n.
\]
Donc \((v_n)\) est géométrique de raison \(3\). Comme
\[
v_0=u_0+2=4,
\]
on obtient :
\[
v_n=4\cdot3^n.
\]
Ainsi :
\[
u_n=v_n-2=4\cdot3^n-2.
\]
Donc :
\[
u_5=4\cdot243-2=970.
\]
\end{solbox}

\begin{solbox}{17}
\[
u_0=7,\quad u_1=4,\quad u_2=3,\quad u_3=4,\quad u_4=7.
\]
Puis :
\[
u_{n+1}-u_n=((n+1)^2-4(n+1)+7)-(n^2-4n+7).
\]
Donc :
\[
u_{n+1}-u_n=2n-3.
\]
Cette quantité est négative pour \(n=0,1\), positive pour \(n\geq2\). La suite décroît jusqu'à \(u_2\), puis croît à partir de \(u_2\).
\end{solbox}

\begin{solbox}{18}
\[
S_0=1,\quad S_1=3,\quad S_2=7,\quad S_3=15.
\]
On conjecture :
\[
S_n=2^{n+1}-1.
\]
C'est une somme géométrique de raison \(2\), donc :
\[
S_n=\frac{1-2^{n+1}}{1-2}=2^{n+1}-1.
\]
On cherche :
\[
2^{n+1}-1\geq1000,
\]
soit
\[
2^{n+1}\geq1001.
\]
Or \(2^9=512\) et \(2^{10}=1024\). Donc \(n+1=10\), soit \(n=9\).
\end{solbox}

\begin{solbox}{19}
\[
u_0=5,\quad u_1=\frac52,\quad u_2=\frac53.
\]
On a :
\[
u_{n+1}=\frac{5}{n+2}<\frac{5}{n+1}=u_n
\]
car \(n+2>n+1>0\). La suite est donc strictement décroissante.  
Elle n'est pas arithmétique car les différences ne sont pas constantes ; elle n'est pas géométrique car les quotients ne sont pas constants.  
On cherche :
\[
\frac{5}{n+1}\leq0{,}1.
\]
Donc :
\[
5\leq0{,}1(n+1),
\]
\[
50\leq n+1,
\]
d'où \(n\geq49\).
\end{solbox}

\begin{solbox}{20}
Les valeurs successives sont :
\[
3,\ 7,\ 15,\ 31,\ 63,\ 127,\ 255.
\]
La suite vérifie :
\[
u_0=3,\qquad u_{n+1}=2u_n+1.
\]
Avec \(v_n=u_n+1\), on obtient :
\[
v_{n+1}=u_{n+1}+1=2u_n+2=2(u_n+1)=2v_n.
\]
Donc \((v_n)\) est géométrique de raison \(2\).
\end{solbox}

\begin{solbox}{21}
On a :
\[
u_1=u_0+r,\quad u_2=u_0+2r,\quad u_3=u_0+3r.
\]
La formule conjecturée est :
\[
u_n=u_0+nr.
\]
En effet, à chaque passage d'un rang au suivant, on ajoute \(r\). Après \(n\) passages à partir de \(u_0\), on a ajouté \(n\) fois \(r\).  
Application :
\[
u_n=-2+5n.
\]
\end{solbox}

\begin{solbox}{22}
\[
u_1=u_0q,\quad u_2=u_0q^2,\quad u_3=u_0q^3.
\]
Donc :
\[
u_n=u_0q^n.
\]
À chaque passage d'un rang au suivant, on multiplie par \(q\). Après \(n\) passages, on a multiplié \(n\) fois par \(q\).  
Avec \(u_0=4\) et \(q=-2\) :
\[
u_n=4(-2)^n.
\]
\end{solbox}

\begin{solbox}{23}
\[
S_n=1+q+q^2+\cdots+q^n.
\]
Alors :
\[
qS_n=q+q^2+\cdots+q^{n+1}.
\]
Donc :
\[
S_n-qS_n=1-q^{n+1}.
\]
Ainsi :
\[
(1-q)S_n=1-q^{n+1}.
\]
Comme \(q\neq1\),
\[
S_n=\frac{1-q^{n+1}}{1-q}.
\]
Avec \(q=\frac12\) et \(n=5\) :
\[
S_5=\frac{1-(1/2)^6}{1-1/2}=2\left(1-\frac1{64}\right)=\frac{63}{32}.
\]
\end{solbox}

\begin{solbox}{24}
Pour une suite arithmétique :
\[
u_{n+1}=u_n+r.
\]
Donc :
\[
u_{n+1}-u_n=r.
\]
Si \(r>0\), cette différence est strictement positive, donc la suite est strictement croissante.  
Si \(r<0\), elle est strictement négative, donc la suite est strictement décroissante.  
Si \(r=0\), la suite est constante.
\end{solbox}

\begin{solbox}{25}
Pour une suite géométrique :
\[
u_{n+1}=qu_n.
\]
Donc :
\[
u_{n+1}-u_n=qu_n-u_n=(q-1)u_n.
\]
Comme \(u_0>0\) et \(q>0\), tous les termes sont positifs.  
Si \(q>1\), alors \(q-1>0\), donc \(u_{n+1}-u_n>0\) : la suite est croissante.  
Si \(0<q<1\), alors \(q-1<0\), donc elle est décroissante.
\end{solbox}

\begin{solbox}{26}
On cherche \(\alpha\) tel que \(v_n=u_n+\alpha\) vérifie \(v_{n+1}=2v_n\).  
Or :
\[
v_{n+1}=u_{n+1}+\alpha=2u_n+3+\alpha.
\]
On veut :
\[
2u_n+3+\alpha=2(u_n+\alpha)=2u_n+2\alpha.
\]
Donc :
\[
3+\alpha=2\alpha,
\]
d'où \(\alpha=3\).  
Ainsi \(v_n=u_n+3\) et
\[
v_{n+1}=2v_n.
\]
Comme \(v_0=4\), on a :
\[
v_n=4\cdot2^n,
\]
donc :
\[
u_n=4\cdot2^n-3.
\]
\end{solbox}

\begin{solbox}{27}
\[
u_0=1,\quad u_1=2,\quad u_2=5,\quad u_3=10,\quad u_4=17.
\]
On reconnaît :
\[
u_n=n^2+1.
\]
Vérification :
\[
u_{n+1}=(n+1)^2+1=n^2+2n+2.
\]
Et :
\[
u_n+2n+1=n^2+1+2n+1=n^2+2n+2.
\]
La formule est compatible avec la récurrence et avec \(u_0=1\).
\end{solbox}

\begin{solbox}{28}
\[
u_{n+1}=\frac{2n+5}{n+3}.
\]
Donc :
\[
u_{n+1}-u_n=\frac{2n+5}{n+3}-\frac{2n+3}{n+2}.
\]
On met au même dénominateur :
\[
u_{n+1}-u_n=\frac{(2n+5)(n+2)-(2n+3)(n+3)}{(n+3)(n+2)}.
\]
Le numérateur vaut :
\[
2n^2+9n+10-(2n^2+9n+9)=1.
\]
Donc :
\[
u_{n+1}-u_n=\frac{1}{(n+3)(n+2)}>0.
\]
La suite est strictement croissante.
\end{solbox}

\begin{solbox}{29}
Pour tout \(n\in\N\),
\[
\frac{2}{n+1}>0,
\]
donc
\[
u_n=3+\frac{2}{n+1}>3.
\]
De plus, comme \(n+1\geq1\),
\[
\frac{2}{n+1}\leq2,
\]
donc
\[
u_n\leq5.
\]
Enfin,
\[
u_{n+1}=3+\frac{2}{n+2}<3+\frac{2}{n+1}=u_n.
\]
La suite est décroissante et reste au-dessus de \(3\).
\end{solbox}

\begin{solbox}{30}
\[
u_n=6\left(\frac13\right)^n.
\]
Donc :
\[
S_n=6\frac{1-\left(\frac13\right)^{n+1}}{1-\frac13}
=6\frac{1-\left(\frac13\right)^{n+1}}{\frac23}
=9\left(1-\left(\frac13\right)^{n+1}\right).
\]
Comme \(\left(\frac13\right)^{n+1}>0\),
\[
S_n<9.
\]
Les termes ajoutés diminuent car la raison \(\frac13\) est comprise entre \(0\) et \(1\).
\end{solbox}

\begin{solbox}{31}
\[
u_0=10,\ u_1=15,\ u_2=20,\ u_3=25,\ u_4=30.
\]
\[
v_0=10,\ v_1=12,\ v_2=14{,}4,\ v_3=17{,}28,\ v_4=20{,}736.
\]
Par essais :
\[
v_{10}\approx61{,}92,\quad u_{10}=60,
\]
donc le dépassement se produit vers \(n=10\).  
La suite \(u\) augmente par addition de \(5\), tandis que \(v\) augmente de \(20\%\) à chaque rang. La croissance en pourcentage s'amplifie car elle s'applique à un terme de plus en plus grand.
\end{solbox}

\begin{solbox}{32}
La perte de \(4\%\) donne une multiplication par \(0{,}96\), puis on ajoute \(200\) :
\[
P_{n+1}=0{,}96P_n+200.
\]
On calcule :
\[
P_1=9800,\quad P_2=9608,\quad P_3=9423{,}68.
\]
La suite n'est ni arithmétique ni géométrique.  
Posons \(Q_n=P_n-5000\). Alors :
\[
Q_{n+1}=P_{n+1}-5000=0{,}96P_n+200-5000
=0{,}96P_n-4800=0{,}96(P_n-5000)=0{,}96Q_n.
\]
Donc \((Q_n)\) est géométrique. Comme \(Q_0=5000\),
\[
Q_n=5000\cdot0{,}96^n.
\]
Ainsi :
\[
P_n=5000+5000\cdot0{,}96^n.
\]
\end{solbox}

\begin{solbox}{33}
\[
D_{n+1}=1{,}01D_n-150.
\]
\[
D_1=4900,\quad D_2=4799,\quad D_3=4696{,}99.
\]
On pose \(E_n=D_n-15000\). Alors :
\[
E_{n+1}=D_{n+1}-15000=1{,}01D_n-150-15000
=1{,}01D_n-15150=1{,}01(D_n-15000)=1{,}01E_n.
\]
Donc :
\[
E_n=E_0\cdot1{,}01^n=(-10000)\cdot1{,}01^n.
\]
Ainsi :
\[
D_n=15000-10000\cdot1{,}01^n.
\]
La dette finit par devenir négative dans le modèle, ce qui signifie qu'elle s'annule avant ce moment. En pratique, on arrête le modèle lorsque la dette atteint \(0\).
\end{solbox}

\begin{solbox}{34}
\[
u_1=7,\quad u_2=\frac{11}{2}=5{,}5,\quad u_3=4{,}75.
\]
Si \(u_n>4\), alors :
\[
u_{n+1}=\frac{u_n+4}{2}>\frac{4+4}{2}=4.
\]
De plus :
\[
u_{n+1}-u_n=\frac{u_n+4}{2}-u_n=\frac{4-u_n}{2}.
\]
Comme \(u_n>4\), cette différence est négative : la suite est décroissante.  
Avec \(v_n=u_n-4\),
\[
v_{n+1}=u_{n+1}-4=\frac{u_n+4}{2}-4=\frac{u_n-4}{2}=\frac12v_n.
\]
Donc :
\[
v_n=6\left(\frac12\right)^n.
\]
Ainsi :
\[
u_n=4+6\left(\frac12\right)^n.
\]
\end{solbox}

\begin{solbox}{35}
Chaque heure, il reste \(70\%\), puis on ajoute \(5\) :
\[
M_{n+1}=0{,}7M_n+5.
\]
Avec \(M_0=5\),
\[
M_1=8{,}5,\quad M_2=10{,}95,\quad M_3=12{,}665.
\]
La valeur stable \(L\) vérifie :
\[
L=0{,}7L+5.
\]
Donc :
\[
0{,}3L=5,\quad L=\frac{50}{3}.
\]
Avec \(V_n=M_n-\frac{50}{3}\),
\[
V_{n+1}=0{,}7V_n.
\]
Donc :
\[
M_n=\frac{50}{3}+\left(5-\frac{50}{3}\right)0{,}7^n.
\]
\end{solbox}

\begin{solbox}{36}
La suite modélise une quantité initiale \(2\), augmentée de \(15\%\) à chaque rang.  
Par essais :
\[
u_{10}=2\cdot1{,}15^{10}\approx8{,}09,
\]
\[
u_{12}=2\cdot1{,}15^{12}\approx10{,}70.
\]
On vérifie \(u_{11}\approx9{,}30\). Le premier rang est donc \(n=12\).  
Algorithme :
\[
\begin{array}{l}
n\leftarrow0,\ u\leftarrow2\\
\text{Tant que }u<10\\
\quad u\leftarrow1{,}15u\\
\quad n\leftarrow n+1\\
\text{Fin Tant que}
\end{array}
\]
Le rang existe car la suite géométrique de raison \(1{,}15>1\) croît sans cesse.
\end{solbox}

\begin{solbox}{37}
Pour une suite géométrique :
\[
u_5=u_2q^3.
\]
Donc :
\[
486=18q^3.
\]
Ainsi :
\[
q^3=27.
\]
Dans \(\R\), on obtient \(q=3\). Donc :
\[
u_n=u_2q^{n-2}=18\cdot3^{n-2}.
\]
\end{solbox}

\begin{solbox}{38}
Pour une suite arithmétique :
\[
u_{12}-u_4=(12-4)r=8r.
\]
Donc :
\[
49-17=8r,
\]
\[
32=8r,
\]
d'où
\[
r=4.
\]
Puis :
\[
u_4=u_0+4r,
\]
donc :
\[
17=u_0+16,
\]
et
\[
u_0=1.
\]
Ainsi :
\[
u_n=1+4n.
\]
\end{solbox}

\begin{solbox}{39}
\[
S_n=3(1+2+\cdots+2^n)=3(2^{n+1}-1).
\]
Pour \(S_n=381\),
\[
3(2^{n+1}-1)=381.
\]
Donc :
\[
2^{n+1}-1=127,
\]
\[
2^{n+1}=128=2^7.
\]
Donc :
\[
n+1=7,\quad n=6.
\]
Pour \(S_n\geq1000\),
\[
3(2^{n+1}-1)\geq1000.
\]
Par essais :
\[
S_8=3(512-1)=1533,
\]
\[
S_7=3(256-1)=765.
\]
Le premier rang est \(n=8\).
\end{solbox}

\begin{solbox}{40}
La relation est :
\[
P_{n+1}=1{,}08P_n-30.
\]
Elle n'est pas arithmétique à cause du facteur \(1{,}08\), ni géométrique à cause du terme \(-30\).  
On cherche \(L\) tel que :
\[
P_{n+1}-L=1{,}08(P_n-L).
\]
Cela revient à imposer :
\[
1{,}08P_n-30-L=1{,}08P_n-1{,}08L.
\]
Donc :
\[
-30-L=-1{,}08L,
\]
\[
30=0{,}08L,
\]
\[
L=375.
\]
Ainsi :
\[
P_n-375=(P_0-375)1{,}08^n.
\]
Donc :
\[
P_n=375+625\cdot1{,}08^n.
\]
Comme \(1{,}08>1\) et \(625>0\), la suite est croissante.
\end{solbox}

\begin{solbox}{41}
\[
u_1=8,\quad u_2=12{,}8,\quad u_3=15{,}68,\quad u_4=17{,}408.
\]
La valeur stable \(L\) vérifie :
\[
L=0{,}6L+8.
\]
Donc :
\[
0{,}4L=8,\quad L=20.
\]
Avec \(v_n=u_n-20\),
\[
v_{n+1}=u_{n+1}-20=0{,}6u_n+8-20=0{,}6u_n-12=0{,}6(u_n-20)=0{,}6v_n.
\]
Donc :
\[
v_n=v_0\cdot0{,}6^n=-20\cdot0{,}6^n.
\]
Ainsi :
\[
u_n=20-20\cdot0{,}6^n.
\]
On cherche :
\[
20-20\cdot0{,}6^n>19,
\]
soit
\[
0{,}6^n<0{,}05.
\]
Par essais, \(0{,}6^5\approx0{,}0778\) et \(0{,}6^6\approx0{,}0467\). Le premier rang est \(n=6\).
\end{solbox}

\begin{solbox}{42}
Si \(C_n\) est le capital juste après le versement du mois \(n\), alors :
\[
C_{n+1}=1{,}005C_n+100.
\]
On peut aussi écrire, après \(n\) mois :
\[
C_n=100(1+1{,}005+1{,}005^2+\cdots+1{,}005^n).
\]
Donc :
\[
C_n=100\frac{1-1{,}005^{n+1}}{1-1{,}005}.
\]
Pour \(24\) mois :
\[
C_{24}=100\frac{1-1{,}005^{25}}{1-1{,}005}\approx2654.
\]
\end{solbox}

\begin{solbox}{43}
\(u_n\) est une croissance arithmétique de \(30\) unités par rang.  
\(v_n\) est une croissance géométrique de \(8\%\) par rang.  
Par essais :
\[
u_{20}=1100,\quad v_{20}\approx932,
\]
\[
u_{30}=1400,\quad v_{30}\approx2013.
\]
Le dépassement a lieu entre \(20\) et \(30\). Un algorithme naturel consiste à initialiser \(n=0\), puis à augmenter \(n\) jusqu'à obtenir \(v_n>u_n\).  
La suite \((v_n)\) est croissante car sa raison \(1{,}08>1\) et son premier terme est positif.
\end{solbox}

\begin{solbox}{44}
\[
u_1=2a+1,\quad u_2=2(2a+1)+1=4a+3,\quad u_3=8a+7.
\]
Avec \(v_n=u_n+1\),
\[
v_{n+1}=u_{n+1}+1=2u_n+2=2(u_n+1)=2v_n.
\]
Donc :
\[
v_n=(a+1)2^n.
\]
Ainsi :
\[
u_n=(a+1)2^n-1.
\]
On cherche :
\[
u_5>100.
\]
Donc :
\[
32(a+1)-1>100.
\]
\[
32(a+1)>101.
\]
\[
a+1>\frac{101}{32}.
\]
\[
a>\frac{69}{32}.
\]
\end{solbox}

\begin{solbox}{45}
\[
u_1=9,\quad u_2=7{,}5,\quad u_3=6{,}75.
\]
Si \(u_n>6\), alors :
\[
u_{n+1}=\frac{u_n+6}{2}>6.
\]
De plus :
\[
u_{n+1}-u_n=\frac{u_n+6}{2}-u_n=\frac{6-u_n}{2}<0.
\]
Donc la suite est décroissante.  
Avec \(v_n=u_n-6\),
\[
v_{n+1}=\frac12v_n.
\]
Comme \(v_0=6\),
\[
v_n=6\left(\frac12\right)^n.
\]
Donc :
\[
u_n=6+6\left(\frac12\right)^n.
\]
On cherche :
\[
6+6\left(\frac12\right)^n<6{,}1.
\]
Donc :
\[
6\left(\frac12\right)^n<0{,}1,
\]
\[
\left(\frac12\right)^n<\frac1{60}.
\]
Or \(2^5=32\) et \(2^6=64\). Le premier rang est \(n=6\).
\end{solbox}

\begin{solbox}{46}
Les premiers termes sont :
\[
3,\ -6,\ 12,\ -24,\ 48,\ -96.
\]
La suite n'est ni croissante ni décroissante car les signes alternent.  
Elle est géométrique car :
\[
u_{n+1}=3(-2)^{n+1}=-2\cdot3(-2)^n=-2u_n.
\]
Sa raison est \(-2\).  
La somme vaut :
\[
S_n=3\frac{1-(-2)^{n+1}}{1-(-2)}
=1-(-2)^{n+1}.
\]
Si \(n\) est pair, \((-2)^n>0\), donc \(u_n>0\). Si \(n\) est impair, \(u_n<0\).
\end{solbox}

\begin{solbox}{47}
La relation est :
\[
A_{n+1}=0{,}95A_n+180,\qquad A_0=2000.
\]
\[
A_1=2080,\quad A_2=2156,\quad A_3=2228{,}2.
\]
La valeur stable \(L\) vérifie :
\[
L=0{,}95L+180.
\]
Donc :
\[
0{,}05L=180,\quad L=3600.
\]
Avec \(V_n=A_n-3600\),
\[
V_{n+1}=0{,}95V_n.
\]
Donc :
\[
A_n=3600+(2000-3600)0{,}95^n
=3600-1600\cdot0{,}95^n.
\]
La suite se rapproche de \(3600\) abonnés.
\end{solbox}

\begin{solbox}{48}
L'épargne hebdomadaire est une suite arithmétique :
\[
u_1=50,\qquad r=5.
\]
Donc :
\[
u_n=50+5(n-1).
\]
La somme sur \(20\) semaines vaut :
\[
S_{20}=\frac{20(u_1+u_{20})}{2}.
\]
Or :
\[
u_{20}=50+5\cdot19=145.
\]
Donc :
\[
S_{20}=10(50+145)=1950.
\]
Pour atteindre \(2000\), il faut une semaine de plus car \(S_{20}=1950\) et \(u_{21}=150\). Donc le nombre minimal est \(21\).
\end{solbox}

\begin{solbox}{49}
On calcule :
\[
u_1=2+3=5,
\]
\[
u_2=5+6=11,
\]
\[
u_3=11+12=23,
\]
\[
u_4=23+24=47.
\]
On a :
\[
u_n=2+3(1+2+2^2+\cdots+2^{n-1})
\]
pour \(n\geq1\).  
La somme vaut :
\[
1+2+\cdots+2^{n-1}=2^n-1.
\]
Donc :
\[
u_n=2+3(2^n-1)=3\cdot2^n-1.
\]
Vérification :
\[
u_0=3-1=2,\quad u_1=6-1=5.
\]
\end{solbox}

\begin{solbox}{50}
Le modèle dit qu'il reste \(90\%\) du stock, puis qu'on ajoute \(12\) unités :
\[
u_{n+1}=0{,}9u_n+12.
\]
\[
u_1=0{,}9\cdot100+12=102,
\]
\[
u_2=0{,}9\cdot102+12=103{,}8,
\]
\[
u_3=0{,}9\cdot103{,}8+12=105{,}42.
\]
La valeur stable \(L\) vérifie :
\[
L=0{,}9L+12.
\]
Donc :
\[
0{,}1L=12,
\]
\[
L=120.
\]
Posons \(v_n=u_n-120\). Alors :
\[
v_{n+1}=u_{n+1}-120=0{,}9u_n+12-120=0{,}9u_n-108=0{,}9(u_n-120)=0{,}9v_n.
\]
Donc \((v_n)\) est géométrique de raison \(0{,}9\). Comme :
\[
v_0=100-120=-20,
\]
on obtient :
\[
v_n=-20\cdot0{,}9^n.
\]
Ainsi :
\[
u_n=120-20\cdot0{,}9^n.
\]
La suite est croissante car \(0{,}9^n\) décroît, donc \(-20\cdot0{,}9^n\) augmente vers \(0\).  
On cherche :
\[
120-20\cdot0{,}9^n>119.
\]
Donc :
\[
20\cdot0{,}9^n<1,
\]
\[
0{,}9^n<0{,}05.
\]
Par essais :
\[
0{,}9^{28}\approx0{,}0523,\qquad 0{,}9^{29}\approx0{,}0471.
\]
Le premier rang est donc \(n=29\).
\end{solbox}

\newpage

%================================================
\section{Fiche de synthèse finale}

\begin{rememberbox}{Formules essentielles}
\[
\boxed{\text{Suite arithmétique : } u_{n+1}=u_n+r}
\]
\[
\boxed{u_n=u_0+nr}
\]
\[
\boxed{\text{Suite géométrique : } u_{n+1}=qu_n}
\]
\[
\boxed{u_n=u_0q^n}
\]
\[
\boxed{1+q+\cdots+q^n=\frac{1-q^{n+1}}{1-q}\quad(q\neq1)}
\]
\end{rememberbox}

\begin{methodbox}{Étudier une suite}
Pour étudier une suite, on peut suivre le plan suivant :
\begin{enumerate}
  \item Identifier le type de définition : explicite ou récurrente.
  \item Calculer quelques premiers termes pour comprendre le comportement.
  \item Chercher si la suite est arithmétique ou géométrique.
  \item Pour les variations, calculer \(u_{n+1}-u_n\).
  \item Pour un modèle en pourcentage, traduire la variation par un coefficient multiplicateur.
  \item Pour une suite du type \(u_{n+1}=au_n+b\), chercher une valeur stable \(L\) telle que \(L=aL+b\).
\end{enumerate}
\end{methodbox}

\begin{warningbox}{Erreurs fréquentes}
\begin{itemize}
  \item Confondre augmentation de \(20\%\) et multiplication par \(20\).
  \item Oublier que \(n\) est un entier naturel.
  \item Dire qu'une suite est croissante après avoir seulement calculé quelques termes.
  \item Utiliser la formule d'une suite arithmétique pour une suite géométrique.
  \item Oublier le nombre exact de termes dans une somme.
\end{itemize}
\end{warningbox}

\end{document}